

\documentclass[12]{amsart}
\usepackage{amsmath, amssymb, amsthm}
\usepackage{geometry}
\usepackage{mathrsfs}
\usepackage{lmodern}
\usepackage{graphicx}
\usepackage{epigraph}
\usepackage{palatino}
\usepackage{titlesec}
\usepackage{setspace}
\usepackage{mdframed}
\usepackage{lipsum}
\usepackage{enumitem}
\usepackage{fancyhdr}
\geometry{margin=1in}
\pagestyle{fancy}
\lhead{The Word of Truth}
\rhead{\thepage}
\fancyfoot[C]{}        % Clear center footer (where default page number appears)
\fancyfoot[L]{}        % Optional: clear left footer
\fancyfoot[R]{}        % Optional: clear right footer
\renewcommand{\headrulewidth}{0.4pt}  % keep or customize header line
\renewcommand{\footrulewidth}{0pt}    
% Define theorem style (optional)
\theoremstyle{plain}
\newtheorem{theorem}{Theorem}[section]
\newtheorem{lemma}[theorem]{Lemma}
\newtheorem{definition}[theorem]{Definition}
\begin{document}

\begin{titlepage}
    \centering
    \vspace*{3cm}
    
{\Huge\bfseries\ The Word of Truth:\\ Aleph, Mem, Tav}
\author{\Large A Simple Revelation for All People}
\date{}
    \vspace{0.5cm}
   \begin{center}
\large\textit{“I am the Aleph and the Tav, the Beginning and the End.” — Revelation 22:13}
\end{center}

    \vfill

    \vspace{1em}
    \begin{minipage}{1\textwidth}
        
This presentation unveils a timeless truth encoded in the Hebrew language—the divine signature of truth itself, written in three sacred letters: \textbf{Aleph (\א)}, \textbf{Mem (\מ)}, and \textbf{Tav (\ת)}. These letters form the Hebrew word \textbf{Emet} (\אמת), meaning \textit{“truth”}.

\medskip

In this revelation, we discover that:
\begin{itemize}
  \item \textbf{Aleph} represents the beginning—Yahuah, the eternal and holy source of all creation.
  \item \textbf{Tav} represents the end—the seal of covenant, completion, and righteous judgment.
  \item At the very center stands \textbf{Mem}—symbolizing the Torah, the living Word, the wellspring of instruction through which all things are established and sustained.
\end{itemize}

\medskip

This structure is no accident. It is the spiritual architecture of reality. The Aleph and the Tav are not only the first and last letters of the Hebrew Alephbet—they are a divine title revealed by Yahusha the Messiah: \textit{“I am the Aleph and the Tav, the Beginning and the End.”} (Revelation 22:13)

\medskip

Through this sacred lens, we see that the \textbf{truth} is not a concept—it is a covenant. It begins in the unity of God, flows through the teaching of His Torah (Mem), and finds its fulfillment in His righteous judgment and loving redemption. Truth is alive. Truth walks among us. Truth is the Word made flesh.

\medskip

May this presentation open the eyes of our hearts to behold the beauty of His Name, the power of His Word, and the unshakable truth that binds heaven and earth: \textbf{Emet}—from Aleph to Tav, held together by the Living Torah.

    \end{minipage}
\end{titlepage}
\clearpage

\begin{center}
\Huge\textbf{\אמת} \\
\large Aleph -- Mem -- Tav
\end{center}

\section*{1. Aleph (א) — The Beginning}
\begin{itemize}
  \item First letter of the Hebrew Alephbet.
  \item Represents unity, source, and divine origin.
  \item Symbolizes the eternal God (Yahuah), the Creator and initiator of all things.
\end{itemize}

\section*{2. Tav (ת) — The End}
\begin{itemize}
  \item Last letter of the Hebrew Alephbet.
  \item Means ``mark'' or ``seal'' — the signature of completion.
  \item Represents fulfillment, covenant, and righteous judgment.
\end{itemize}

\section*{3. Mem (מ) — The Middle, the Heart}
\begin{itemize}
  \item Middle letter of the word \textbf{אמת} (Emet, “truth”).
  \item Symbolizes \textbf{Torah} — the instructions (teachings) of God.
  \item Mem represents water, flow, and multiplication — God's Word bringing life to the world.
  \item Mem connects the Beginning (Aleph) to the End (Tav).
\end{itemize}

\section*{The Message of TRUTH (אמת)}
\begin{itemize}
  \item Truth is not static—it begins with God (Aleph), is revealed through His Word (Mem), and brings all things to fulfillment (Tav).
  \item Yeshua (Jesus) said, “I am the Way, the \textbf{Truth}, and the Life.” (John 14:6) — He is the Aleph, the Mem, and the Tav.
  \item The Torah (God's instruction) is the center of truth, just as Yeshua is the Living Word made flesh (John 1:14).
\end{itemize}

\section*{The Word \textbf{Amen} (אמן)}
\begin{itemize}
  \item Also begins with Aleph, contains Mem, and ends in Nun, which symbolizes continuity and life.
  \item When we say “Amen,” we affirm: “Let this truth be established forever.”
\end{itemize}

\section*{Conclusion}
\begin{quote}
The Hebrew word for truth—\textbf{אמת}—is a divine message encoded in the language of Scripture:

\begin{center}
\textbf{Aleph (God) + Mem (Torah) + Tav (Son) = Truth}
\end{center}

This message is a blessing for all people: God is the source, His Word is life, and He brings all things to completion in love and righteousness.
\end{quote}

\vspace{1em}
\begin{center}
\textit{“Sanctify them by Your truth; Your Word is truth.” — John 17:17}
\end{center}
\end{document}
